% ARCHIVO GENERADO por scripts/generar_capitulos.py. No editar a mano.
\section{Historias de usuario de la iteración 3}
Aprobadas por el equipo el 2026-11-20 (versión 1), con los cuatro puntos de revisión de la última subsección.

\subsection{Resumen}

\begin{longtable}{>{\raggedright\arraybackslash}p{1.5cm}>{\raggedright\arraybackslash}p{7.0cm}>{\raggedright\arraybackslash}p{2.6cm}>{\raggedright\arraybackslash}p{1.5cm}>{\raggedright\arraybackslash}p{1.8cm}}
\toprule
\textbf{ID} & \textbf{Título} & \textbf{Actor} & \textbf{RF} & \textbf{Prioridad}\\
\midrule\endhead
HU-024 & Registrar el pago de una mensualidad & Operador & RF-201 & Must\\
HU-025 & Consultar la vigencia de una mensualidad & Operador & RF-202 & Must\\
HU-026 & Consultar el historial de pagos e imprimir la factura & Operador & RF-203 & Must\\
HU-027 & Gestionar las tarifas & Administrador & RF-204 & Must\\
HU-028 & Inactivar con motivo y reactivar al pagar & Operador & RF-205 & Must\\
HU-029 & Corregir un pago mal registrado & Operador, Administrador & RF-206 & Must\\
HU-030 & Ver las mensualidades por vencer & Operador & RF-207 & Should\\
HU-031 & Gestionar los descuentos & Administrador & RF-208 & Must\\
HU-032 & Consultar los ingresos por mes & Administrador & RF-209 & Could\\
\bottomrule
\end{longtable}




\subsection{HU-024 — Registrar el pago de una mensualidad}

\textbf{Historia:} Como \textbf{operador}, quiero registrar el pago de uno o varios meses de un vehículo, con el valor calculado por el sistema, para que el vehículo quede al día y el cliente reciba su factura.


\textbf{Requerimientos:} RF-201 · \textbf{Reglas:} RN-03, RN-17, RN-38 a RN-43, RN-46, RN-47 · \textbf{Prioridad:} Must


\begin{lstlisting}[style=gherkin]
Escenario: Primer pago de un vehículo
  Dado que "ABC123" no tiene pagos
  Cuando el operador registra 1 mes en efectivo el 20/11/2026
  Entonces el pago cubre del 20/11/2026 al 20/12/2026 por 150 000 COP
  Y el día de cobro del vehículo queda en 20
  Y la bitácora muestra que el operador registró el pago

Escenario: Antes de confirmar, el sistema muestra el resumen
  Cuando el operador escribe la placa "ABC123", 1 mes y el medio de pago
  Entonces el sistema muestra el propietario, la tarifa, el periodo que cubrirá y el valor
  Y el pago solo se guarda cuando el operador confirma

Escenario: Pagar tarde no mueve el día de cobro
  Dado que la mensualidad de "ABC123" venció el 18/11/2026
  Cuando el operador registra 1 mes el 20/11/2026
  Entonces el pago cubre del 18/11/2026 al 18/12/2026

Escenario: Fin de mes
  Dado que el día de cobro de un vehículo es 31 y su mensualidad vence el 31/01/2027
  Cuando se registran 2 meses
  Entonces el pago cubre hasta el 31/03/2027, pasando por el 28/02/2027

Escenario: Varios meses atrasados
  Dado que la mensualidad de "ABC123" venció el 18/08/2026 y el vehículo no ha vuelto
  Cuando el operador busca la placa para registrar el pago
  Entonces el sistema indica que debe 4 meses para quedar al día
  Y si el operador registra solo 1 mes, el pago cubre del 18/08/2026 al 18/09/2026
  Y el sistema advierte que el vehículo sigue con la mensualidad vencida

Escenario: Varios meses en un solo pago
  Cuando el operador registra 3 meses para "ABC123"
  Entonces el valor es 450 000 COP
  Y el número de meses debe estar entre 1 y 12

Escenario: Aplicar un descuento
  Dado un descuento habilitado "Cliente frecuente" de 10 %
  Cuando el operador registra 3 meses con ese descuento
  Entonces el valor es 405 000 COP
  Y el pago guarda el valor sin descuento, el descuento y el valor final

Escenario: El valor no se escribe a mano
  Cuando el operador registra un pago
  Entonces el valor lo calcula el sistema y no se puede editar

Escenario: Medio de pago obligatorio
  Cuando el operador no elige efectivo, transferencia o tarjeta
  Entonces el sistema no guarda e indica el campo con error

Escenario: Vehículo dentro con la mensualidad vencida
  Dado que "ABC123" está dentro con la mensualidad vencida hace 3 días
  Cuando el operador registra el pago y el vehículo queda al día
  Entonces el vehículo sale del proceso de retiro y no se envían más avisos

Escenario: Confirmación al propietario
  Cuando se registra el pago
  Entonces se envía un correo simulado a "laura@correo.co" con el periodo y el valor
  Y el operador puede abrir la factura para imprimirla
\end{lstlisting}



\subsection{HU-025 — Consultar la vigencia de una mensualidad}

\textbf{Historia:} Como \textbf{operador}, quiero ver si la mensualidad de un vehículo está al día y hasta cuándo, para saber si puede entrar y cuánto debe.


\textbf{Requerimientos:} RF-202 · \textbf{Reglas:} RN-03, RN-38, RN-39 · \textbf{Prioridad:} Must


\begin{lstlisting}[style=gherkin]
Escenario: Al día
  Dado que la mensualidad de "ABC123" vence el 18/12/2026
  Cuando el operador consulta el vehículo
  Entonces ve "Al día", el vencimiento 18/12/2026, el día de cobro 18 y el último pago

Escenario: Vence hoy
  Dado que la mensualidad vence el 20/11/2026
  Cuando el operador consulta el vehículo
  Entonces ve "Vence hoy" y el vehículo todavía puede entrar

Escenario: Vencida
  Dado que la mensualidad venció el 15/11/2026
  Cuando el operador consulta el vehículo
  Entonces ve "Vencida hace 5 días" y que debe 1 mes

Escenario: Sin pagos
  Dado un vehículo registrado que nunca ha pagado
  Cuando el operador lo consulta
  Entonces ve "Sin pagos" y el vehículo no puede entrar
\end{lstlisting}



\subsection{HU-026 — Consultar el historial de pagos e imprimir la factura}

\textbf{Historia:} Como \textbf{operador}, quiero ver todos los pagos de un vehículo o de un propietario y abrir la factura de cada uno, para responder las preguntas del cliente.


\textbf{Requerimientos:} RF-203 · \textbf{Reglas:} RN-44, RN-47 · \textbf{Prioridad:} Must


\begin{lstlisting}[style=gherkin]
Escenario: Historial de un vehículo
  Cuando el operador abre el historial de pagos de "ABC123"
  Entonces ve cada pago, del más reciente al más antiguo, con periodo, meses, valor, descuento, medio, fecha y quién lo registró

Escenario: Historial de un propietario
  Dado que "Laura Méndez" tiene dos vehículos
  Cuando el operador abre su historial de pagos
  Entonces ve los pagos de ambos vehículos, con la placa de cada uno

Escenario: Pago corregido en el historial
  Dado un pago que fue corregido
  Cuando el operador abre el historial
  Entonces ese pago aparece marcado como "Corregido", con el motivo y el pago que lo reemplaza

Escenario: Imprimir la factura
  Cuando el operador abre la factura de un pago
  Entonces ve una página lista para imprimir con número consecutivo, fecha, datos del parqueadero, propietario y cédula, vehículo, periodo, meses, tarifa, descuento, total, medio de pago y quién lo registró
  Y la página lleva la marca "Simulación – sin validez fiscal"
  Y el navegador la puede guardar como PDF
\end{lstlisting}



\subsection{HU-027 — Gestionar las tarifas}

\textbf{Historia:} Como \textbf{administrador}, quiero registrar una tarifa nueva por tipo de vehículo con la fecha desde la que rige, para cambiar los precios sin afectar los pagos anteriores.


\textbf{Requerimientos:} RF-204 · \textbf{Reglas:} RN-41, RN-45, DM-04 · \textbf{Prioridad:} Must


\begin{lstlisting}[style=gherkin]
Escenario: Tarifa nueva con fecha futura
  Dado que la tarifa de carro es 150 000 COP
  Cuando el administrador registra 160 000 COP para carro vigente desde el 01/01/2027
  Entonces los pagos registrados antes del 01/01/2027 usan 150 000 COP
  Y los registrados desde el 01/01/2027 usan 160 000 COP

Escenario: Los pagos anteriores no cambian
  Cuando el administrador registra una tarifa nueva
  Entonces los pagos ya registrados conservan su valor

Escenario: Las tarifas no se editan
  Cuando el administrador consulta las tarifas
  Entonces ve la tarifa actual y el historial por tipo
  Pero no puede editar ni borrar una tarifa registrada

Escenario: Fecha repetida
  Dado una tarifa de moto vigente desde el 01/01/2027
  Cuando el administrador registra otra tarifa de moto con la misma fecha
  Entonces el sistema la rechaza con "Ya existe una tarifa de moto desde esa fecha"

Escenario: Solo el administrador
  Cuando un operador intenta registrar una tarifa
  Entonces el sistema responde "acceso denegado"
\end{lstlisting}



\subsection{HU-028 — Inactivar con motivo y reactivar al pagar}

\textbf{Historia:} Como \textbf{operador}, quiero indicar por qué inactivo un vehículo o un propietario, para que el sistema lo reactive solo cuando el motivo era la falta de pago.


\textbf{Requerimientos:} RF-205 · \textbf{Reglas:} RN-46 · \textbf{Prioridad:} Must


\begin{lstlisting}[style=gherkin]
Escenario: Motivo obligatorio
  Cuando el operador inactiva un vehículo
  Entonces debe elegir "Falta de pago" u "Otro"
  Y si elige "Otro", debe escribir la descripción

Escenario: El motivo se ve en la ficha
  Dado un vehículo inactivo por "Falta de pago"
  Cuando el operador abre su ficha
  Entonces ve el motivo, quién lo inactivó y la fecha

Escenario: Reactivación al pagar
  Dado que "ABC123" está inactivo por "Falta de pago"
  Cuando el operador registra un pago que lo deja al día
  Entonces el vehículo queda activo
  Y la bitácora registra la reactivación automática

Escenario: Propietario inactivo por falta de pago
  Dado que "Laura Méndez" está inactiva por "Falta de pago"
  Cuando se registra un pago de uno de sus vehículos
  Entonces la propietaria también queda activa

Escenario: Inactivo por otro motivo
  Dado que "ABC123" está inactivo por "Otro: vehículo vendido"
  Cuando el operador registra un pago
  Entonces el vehículo sigue inactivo
  Y el sistema advierte que debe activarse a mano

Escenario: El sistema no inactiva por su cuenta
  Dado que la mensualidad de un vehículo lleva meses vencida
  Entonces el vehículo sigue activo hasta que un usuario lo inactive
\end{lstlisting}



\subsection{HU-029 — Corregir un pago mal registrado}

\textbf{Historia:} Como \textbf{operador}, quiero corregir un pago que registré mal, sin borrarlo, para que la vigencia y los ingresos queden bien y quede constancia del error.


\textbf{Requerimientos:} RF-206 · \textbf{Reglas:} RN-07, RN-44 · \textbf{Prioridad:} Must


\begin{lstlisting}[style=gherkin]
Escenario: Corrección dentro de las 24 horas
  Dado un pago de 1 mes registrado por el operador hace 2 horas, que debía ser de 2 meses
  Cuando el operador lo corrige con el motivo "Eran 2 meses" y registra los datos correctos
  Entonces el pago original queda marcado como "Corregido" con el motivo
  Y se registra un pago nuevo de 2 meses que lo reemplaza
  Y la vigencia se calcula solo con el pago nuevo
  Y la bitácora registra la corrección

Escenario: Motivo obligatorio
  Cuando el operador intenta corregir un pago sin escribir el motivo
  Entonces el sistema no guarda e indica el campo con error

Escenario: El operador después de 24 horas
  Dado un pago registrado hace 25 horas
  Cuando un operador intenta corregirlo
  Entonces el sistema se lo impide con "Este pago solo lo puede corregir el administrador"

Escenario: El administrador en cualquier momento
  Dado un pago registrado hace 10 días
  Cuando el administrador lo corrige con un motivo
  Entonces la corrección se guarda

Escenario: Un pago corregido no se vuelve a corregir
  Dado un pago marcado como "Corregido"
  Entonces no muestra la opción de corregir; se corrige el pago que lo reemplazó

Escenario: Factura del pago corregido
  Cuando se abre la factura de un pago corregido
  Entonces la página muestra "Anulada – reemplazada por la factura N.º" y el número de la nueva
\end{lstlisting}



\subsection{HU-030 — Ver las mensualidades por vencer}

\textbf{Historia:} Como \textbf{operador}, quiero ver qué mensualidades vencen en los próximos 5 días, para recordarles a los clientes que deben pagar.


\textbf{Requerimientos:} RF-207 · \textbf{Prioridad:} Should


\begin{lstlisting}[style=gherkin]
Escenario: Lista de por vencer
  Dado que "ABC123" vence el 23/11/2026 y "XYZ789" vence el 30/11/2026
  Cuando el operador abre "Por vencer"
  Entonces ve "ABC123" con su propietario, su celular y la fecha de vencimiento
  Pero no ve "XYZ789"

Escenario: Incluye las que vencen hoy
  Dado un vehículo que vence el 20/11/2026
  Cuando el operador abre "Por vencer"
  Entonces lo ve marcado como "Vence hoy"

Escenario: Solo vehículos activos
  Dado un vehículo inactivo que vence el 22/11/2026
  Cuando el operador abre "Por vencer"
  Entonces no lo ve
\end{lstlisting}



\subsection{HU-031 — Gestionar los descuentos}

\textbf{Historia:} Como \textbf{administrador}, quiero crear descuentos con una descripción y un porcentaje y habilitarlos o deshabilitarlos, para que los operadores puedan aplicarlos al cobrar.


\textbf{Requerimientos:} RF-208 · \textbf{Reglas:} RN-42 · \textbf{Prioridad:} Must


\begin{lstlisting}[style=gherkin]
Escenario: Crear un descuento
  Cuando el administrador crea "Cliente frecuente" con 10 %
  Entonces el descuento queda habilitado
  Y la bitácora muestra que el administrador lo creó

Escenario: Porcentaje inválido
  Cuando el administrador escribe 0 % o 120 %
  Entonces el sistema lo rechaza con "El porcentaje debe estar entre 1 y 100"

Escenario: Deshabilitar
  Dado el descuento "Cliente frecuente" habilitado
  Cuando el administrador lo deshabilita
  Entonces el operador ya no lo ve al registrar un pago
  Y los pagos que ya lo usaron no cambian

Escenario: No se edita
  Cuando el administrador consulta un descuento
  Entonces no puede cambiar su porcentaje; para otro valor crea un descuento nuevo

Escenario: El operador solo aplica
  Cuando un operador intenta crear o deshabilitar un descuento
  Entonces el sistema responde "acceso denegado"
  Pero al registrar un pago puede elegir uno habilitado
\end{lstlisting}



\subsection{HU-032 — Consultar los ingresos por mes}

\textbf{Historia:} Como \textbf{administrador}, quiero ver cuánto se recaudó cada mes, para conocer los ingresos del parqueadero.


\textbf{Requerimientos:} RF-209 · \textbf{Reglas:} RN-44 · \textbf{Prioridad:} Could


\begin{lstlisting}[style=gherkin]
Escenario: Ingresos de un mes
  Dado que en noviembre de 2026 se registraron pagos por 1 200 000 COP
  Cuando el administrador consulta los ingresos de noviembre de 2026
  Entonces ve el total y el detalle por tipo de vehículo y por medio de pago

Escenario: Sin pagos corregidos
  Dado un pago de 150 000 COP corregido y reemplazado por uno de 300 000 COP
  Cuando el administrador consulta los ingresos del mes
  Entonces solo cuenta los 300 000 COP

Escenario: Se cuenta por fecha de pago
  Dado un pago registrado el 20/11/2026 que cubre 3 meses
  Cuando el administrador consulta los ingresos
  Entonces todo el valor se cuenta en noviembre de 2026

Escenario: Solo el administrador
  Cuando un operador intenta abrir el informe de ingresos
  Entonces el sistema responde "acceso denegado"
\end{lstlisting}



\subsection{Matriz de trazabilidad RF $\rightarrow$ HU (iteración 3)}

\begin{longtable}{>{\raggedright\arraybackslash}p{11.1cm}>{\raggedright\arraybackslash}p{3.3cm}}
\toprule
\textbf{RF} & \textbf{HU}\\
\midrule\endhead
RF-201 Registrar pago & HU-024\\
RF-202 Vigencia & HU-025\\
RF-203 Historial de pagos & HU-026\\
RF-204 Tarifas & HU-027\\
RF-205 Motivo de inactivación y reactivación & HU-028\\
RF-206 Corregir un pago & HU-029\\
RF-207 Por vencer & HU-030\\
RF-208 Descuentos & HU-031\\
RF-209 Ingresos por mes & HU-032\\
\bottomrule
\end{longtable}


Los 9 RF de la iteración 3 tienen una historia.


\subsection{Puntos revisados (2026-11-20)}

Estos detalles no estaban en los requerimientos; el equipo los aprobó tal como se propusieron:


\begin{enumerate}
  \item \textbf{HU-024:} se permite pagar menos meses de los que se deben; el pago se guarda y el sistema advierte que el vehículo sigue vencido.
  \item \textbf{HU-028:} pagar un vehículo reactiva también al propietario si estaba inactivo por falta de pago.
  \item \textbf{HU-029:} la factura del pago corregido queda como ``Anulada'' y remite a la nueva.
  \item \textbf{HU-032:} los ingresos se cuentan en el mes en que se registró el pago, no repartidos entre los meses que cubre.
\end{enumerate}